\section{おわりに}
% 本研究ではアンサンブルベースの欠損補完方法の代表的手法であるmissForestとの比較において特徴量選択を導入した．ベンチマークデータによる性能評価により改善が確認できた．

本研究では，アンサンブルベースの欠損補完手法の代表例である missForest を基盤とし，複数の特徴選択基準を導入した欠損値補完手法を提案した．提案手法では，補完対象となる特徴量ごとに有効な説明変数の部分集合を構築し，各部分集合に基づいて学習したランダムフォレストを OOB 指標により統合することで，特徴選択の多様性と補完性能の安定化を図った．
ベンチマークデータおよび OpenML データセットを用いた実験により，提案手法は missForest や kNN 補完と比較して，多くの条件下で同等以上の補完精度を示すことを確認した．特に，中程度以上の欠損率や特徴量間の相関構造を持つデータにおいて，特徴選択とアンサンブル統合の効果が有効に機能することが示唆された．
一方で，本研究では連続変数のみを対象としており，カテゴリカル変数を含むデータへの適用や，欠損メカニズムが MCAR 以外の場合の検討は今後の課題である．今後は，より複雑な欠損構造や実データへの適用を通じて，提案手法の汎用性と実用性について検証を進めていく予定である．